\subsection{The type primes and the ramified auxiliary primes}
\label{subsec:special-prime-lengths}
\hypertarget{source-special-prime-lengths}{}

The unramified models leave two different sources of exceptional local behavior. A coefficient prime can equal one of the levels imposing the tame type, or it can ramify in one of the auxiliary quadratic fields. We retain the original row in both calculations.

Keep the cut, its good eigenspaces, and its attached
integral Galois representations fixed. Write \(R=r+t\)
for the full rank and normalize the length at a coefficient
prime \(\ell\) by \(v_\ell(\ell)=1\). Both calculations retain the global terms of the numerator: the auxiliary signature length \(s_\pi\), the strict cohomology sum \(V_\ell\), and the class number \(h_E\).

\subsubsection*{The type-prime hypotheses}

At a type prime the local lifting problem retains both tame characters and labeled Hodge jumps. Complete splitting allows them to be specified at one embedding at a time.

Fix a twisted localization at \(\ell=g_i\), with \(i\in\{0,1\}\), and suppose
it has a good point. All completions of \(E\) at \(\ell\)
have degree one. Each completion determines the unique
labeled embedding of \(E\) inducing that place. The CM
type gives a sign at each such embedding, reversed at
the conjugate place. In the family at this localization
the jump set has three parts: the cyclotomic jumps
\(0,\ldots,t-1\), the jumps within \(\eta+1\) of
\(c=(t-1)/2\), and the jumps in the signed outer range.
The negative sign reflects about the central point \(c=(t-1)/2\):
\[
 a\longmapsto 2c-a=t-1-a.
\]

This preserves the central interval and cyclotomic
set and reverses the signed range about \(c\).
The bound on absolute jumps is \(2b\), so this gives
a common set of at most \(R\) distinct integer jumps
for all members of the family at the place. The jumps
avoid \(-1,t\) by the window conditions.

Take \(\eta=\lfloor\eta_c\rfloor\),
\[
 s_T=\lfloor3b\rfloor+2
\]
and \(S_T=0\) in the signed outer range. Then
\(\eta+t+1<s_T\) since
\[
 \eta_c+S_c+t+2\le b.
\]

In each sign the distance of an outer jump in its
specified direction from \(c\) is at most \(s_T-c\).
For the size requirement of the paired-character
argument, put \(y=kb\), where \(k\ge5\) is the ambient rank parameter. For every fixed \(C_T\le b\)
we have \(s_T\le5b\le y\), \(R+1\le3y\) and \(y\ge5\).
Consequently
\[
 C_T\bigl((s_T+S_T)(R+1)^2+(R+1)^3\bigr)
 \le36y^4<y^8\le g_i.
\]
Both type primes lie in the same interval \([Y,2Y]\).
In particular
\[
 g_{1-i}\le2g_i\le g_i^2.
\]

These bounds and the sign calculation supply the original-field and size conditions on the cut family.
\label{cor:special-prime-lengths-1}

In the semisimplification at a retained point the
rank-\(r\) part and the \(t\) cyclotomic characters
carry the specified common inertia and weight labels.
The inertia comparison is transported along the actual
norm choice, as in Section~\ref{subsec:selected-cut}.
The tame generators and wild triviality at the auxiliary
and signature places come from the local type there.
The nontrivial inertia value at each signature place
is part of the marked data. The kill comparison gives
finiteness of the signature length and the strict
cohomology bound gives finite \(V_\ell\) at all the
required positive exponents. Thus the paired-character estimate may be applied on this cut whenever a good point occurs.
If there are no good points, the projected and intersected
lattices are both zero and the contribution is zero.
\label{cor:special-prime-lengths-2}

\subsubsection*{Lifting on a common tame type}

The zero-monodromy and jump conditions at a type place control each paired character. For the numerator we need a uniform bound on the span of the corresponding pushouts, with one exponent for all characters of the common tame type.

\begin{lemma}[A normalized bound on the tame pushout span]\label{lem:common-tame-pushout}
At a type place of degree one let the source and target of the
surjection of stable lattices be rational subquotients of sums from
one potentially crystalline family with zero monodromy. Use a
Frobenius-stable tame character set of size at most \(N=R+1\)
and at most \(N+1\) jumps in \([-u,H]\), each set containing zero.
Assume \(1\le u<H\) and \(C_T(HN^2+N^3)<\ell\) in the integral
tame lifting comparison. If the finite lifting allowance is \(\ell^e\), let \(D(T_0)\) denote the image of the pushout span modulo the finite condition for the target lattice \(T_0\). Its normalized length satisfies
\[
 \operatorname{len}D(T_0)\le2e\operatorname{rk}T_0.
\] The bound includes the torsion in
integral finite cohomology.
\end{lemma}

\begin{proof}
\proofstep{Put the source and target on one tame extension and bound the periods of character differences.}
At a completion of degree one, zero monodromy and
the tame inertia polynomial give a common finite
extension on which the family is crystalline.
One can form it using the kernel of the finite
inertia characters and an unramified extension for
the Frobenius action on them. The Artin character
of the cyclotomic block is included before taking
its inverse twist. Thus the common character set
for the twisted family contains the trivial
character, has at most \(N=R+1\) elements, and
is Frobenius stable. The dual uses the inverse
set and the same cardinality bound. The
coefficient comparison uses a tame exponent
dividing \(\ell^{(\ell-1)!}-1\).

Write the tame characters additively, so Frobenius
acts by multiplication by \(\ell\). If two labels have periods \(a,b\), Frobenius to the power \(\operatorname{lcm}(a,b)\) fixes both. Thus
\[
 \operatorname{period}(\chi-\chi')\mid\operatorname{lcm}(a,b).
\]
In distinct orbits with
lengths \(a,b\), we have \(a+b\le N\) and hence
\[
 \operatorname{lcm}(a,b)\le ab\le\tfrac14(a+b)^2\le N^2/4.
\]
In one orbit, the difference has period dividing
the common period \(a\). A nonzero orbit omits the
fixed point zero, so \(a\le N-1\le N^2/4\).
The zero difference has period one. The maximum \(A\)
over differences thus satisfies \(1\le A\le
\lfloor N^2/4\rfloor\), including differences with
zero.

The integral tame period comparison is used for a
surjection of stable lattices whose rational source
and target are subquotients of sums of the common
potentially crystalline family. It uses the
Frobenius-stable tame character set of size at most
\(N\), including zero, and a jump set with at most
\(N+1\) entries in \([-u,H]\), including zero,
with \(1\le u<H\).
The comparison uses both the bound on \(\ell\)
\[
                  C_T(HN^2+N^3)<\ell
\]
and the degree-one base. It gives the filtered
integral lift after multiplying by
\(\ell^{C_T(NA+u(N+A))}\). Here is the exponent
calculation. The Breuil--Kisin comparison with tame
descent gives an exact triple for the kernel,
source and target lattices.

Write \(A\) for the
maximum period above. The integral Frobenius
coordinates represent the finite target class
after loss \(e_W=2(N-1)A\) by an element in the
kernel of monodromy. A further factor \(\ell^u\)
puts the element, up to a Frobenius coboundary,
in the sum of the two integral coefficient
submodules used for lifting. The exact-triple
comparison lifts it with zero monodromy and loss
\[
             k_{\mathrm{lift}}
               =u(A+1)+e_W+u+(N+1)u+e_W .
\]
The integral cocycle comparison on the source
needs one more loss \(u(A+1)\). Its image on the
target is the prescribed multiple of the class
modulo an integral coboundary, by functoriality
of the Frobenius-coordinate calculation.
The sum of the four losses is
\[
                   6(N-1)A+u(2A+N+5)
                   \le 7\{NA+u(N+A)\}.
\]

In particular \(C_T\ge40\) suffices. An integral
linear retraction of the enlarged coefficients
preserves the cochain and lattice identities and
returns the lift to the original coefficients.
In our use \(u=\eta+t+1\), \(H=s_T+S_T\),
with the Tate dual on the negative side.
\label{cor:special-prime-lengths-3}

For the local cocycle length, take a finite
collection of cocycles from the cut. Write the
underlying cocycles as images of \(gE-E\), with \(E\)
a vector in a finite sum of the twisted family and
\(gE-E\) in a stable finitely generated lattice. Their
representations can be combined in one such sum
and one stable rational subrepresentation \(X\) of it.
The stable source lattices \(L\) in \(X\)
give finite cocycles \(y\) with values in \(L\),
since \(gE-E\) is a rational coboundary.

The residual maps in the presentation give equivariant
maps
\[
 \phi:L\to U/T_0,
\]
where \(T_0\) is the
target lattice in its rational representation \(U\).
The corresponding classes are represented by \(\phi y\).
The rational subrepresentation comparison puts \(X,U\),
or their Tate duals for the dual lift, in the same sums
of the common potential family. The period estimate
can therefore be applied to a surjection of lattices
between their direct sums.

Write \(K\) for the local base and \(F\) for the coefficient
field. To see the integral comparison, let \(\psi:X\to U\)
be a coefficient-field linear lift of \(\phi\) on \(L\).
Choose it on a basis of the free lattice \(L\). Its
equivariance defect takes values in \(T_0\) on \(L\).
The lattice
\[
 L_\psi=\{(a,b)\in L\oplus U:b-\psi a\in T_0\}
\]
in \(X\oplus U\) is stable for the product action and
surjects onto \(L\). If the lift with loss \(\ell^e\)
is available on this side, lift \(\ell^e[y]\) to an
integral finite cocycle of \(L_\psi\). Its second
coordinate reduces to \(\ell^e\phi y\) modulo a
coboundary in \(U/T_0\) and belongs to the finite image.

For the dual lift, use \(X'=X^\vee(1)\),
\(U'=U^\vee(1)\), and the dual lattices \(L'\) and \(T_0'\)
with respect to the Tate pairings. The transpose
\(\psi^{\rm t}:U'\to X'\) has equivariance defect in \(L'\)
on \(T_0'\), as follows by evaluating the defect of \(\psi\)
on \(L\). Use the lattice
\[
 \{(a',b')\in T_0'\oplus X':b'-\psi^{\rm t}a'\in L'\}
\]
and its projection onto \(T_0'\).

For any
\(d\in H^1_f(K,T_0')\), lift \(\ell^e d\) to a finite
cocycle \((a',b')\) in this lattice. Take the cup product with the dual cocycle first.
For \(g,h\in G_K\) its value in \((F/O)(1)\) is computed by
the following congruences in \(F(1)\):
\begin{equation}\label{eq:local-dual-lift-cup}
 \begin{aligned}
 \langle a'(g),g\psi(y(h))\rangle
 &\equiv\langle a'(g),\psi(gy(h))\rangle\\
 &=\langle\psi^{\rm t}a'(g),gy(h)\rangle
 \equiv\langle b'(g),gy(h)\rangle \pmod{O(1)} .
 \end{aligned}
\end{equation}
The first difference in \eqref{eq:local-dual-lift-cup} pairs \(a'(g)\in T_0'\) with the
equivariance defect in \(T_0\). The second pairs an element of
\(L'\) with \(gy(h)\in L\). Thus the cup cochain for
\(a'\) and \(\phi y\) is the reduction of the rational cup
cochain for \(b'\) and \(y\). The latter rational
cup product vanishes in cohomology by Bloch--Kato
orthogonality.

Integral local Tate duality identifies the
finite image in \(H^1(K,U/T_0)\) with the annihilator of
\(H^1_f(K,T_0')\). It follows that \(\ell^e[\phi y]\)
is in the finite image also in this comparison.

\proofstep{Apply finite-condition duality and count the free and torsion generators.}
The common tame type has supplied an exponent for the local lifts. To convert that exponent into a length for the span, one must count both the free generators of the finite condition and its torsion.

For the length, write
\[
 M=H^1_f(K,T_0')
\]
for the
inverse image of the rational finite subspace. It is
finite over \(O\) by local cohomology finiteness.
The quotient of the pushout span by its intersection
with the finite image injects by the cup pairing into
\(\operatorname{Hom}_O(M/\ell^e M,F/O)\).
The torsion of \(M\) has at most \(\operatorname{rk}T_0\)
generators. Indeed the coefficient exact sequence identifies this torsion module as
\[
 H^0(K,U'/T_0')/\operatorname{im}H^0(K,U').
\]
The latter rational subspace meets \(T_0'\)
in a saturated sublattice, so after quotienting by
its image the ambient divisible module is still a direct
sum of at most \(\operatorname{rk}T_0\) copies of \(F/O\).

The rank of \(M\) is
\(\dim_F D_{\mathrm{dR}}(U')/\mathrm{Fil}^0+
\dim_F H^0(K,U')\) by the finite dimension formula.
The invariant vectors inject into \(\mathrm{Fil}^0\),
so the rank is at most \([K:\mathbb Q_\ell]\operatorname{rk}T_0\).
Put \(r_0=\operatorname{rk} T_0\) and
\(e_O=\operatorname{ord}_O(\ell)\). A free generator of \(M\)
contributes \(ee_O\) to the quotient length, and a torsion
summand \(O/(\varpi^a)\) contributes \(\min(a,ee_O)\).
The two generator bounds therefore give
\[
 \operatorname{length}_O(M/\ell^e M)
 \le ee_O\bigl([K:\mathbb Q_\ell]r_0+r_0\bigr)
 =2ee_O r_0\quad(K=\mathbb Q_\ell).
\]
Its dual has the same length. Dividing by \(e_O\)
gives \(2e\,\operatorname{rk}T_0\) in the normalized
comparison. The bound holds for every finite span with
the same exponent and rank, hence for the full span.
\cite[Proposition 3.8 and Corollary 3.8.4]{bloch-kato}
\label{cor:special-prime-lengths-4}
\end{proof}

\subsubsection*{Passing the paired-character bound to the numerator}

The estimate on a span is made at one local character. To put it in the global numerator, sum it first with the actual good multiplicities and then with the normalized weights of the places above the type prime.

\begin{proposition}[Sum of type-prime lengths on the selected row]\label{prop:type-prime-row-length}
At \(\ell=g_i\), write \(M_\ell=C_*x^{C_*}\eta_c+
2\ell^{-2}+V_\ell+v_\ell(h_E)\), with the fixed \(C_*\) below.
The good multiplicities on the cut have
\[
 \sum_{w\mid\ell,\ w\text{ twisted}}
       e_w\omega_w\eta_{j,w}\le d_j(\log\ell)M_\ell.
\]
The term from \(C_*x^{C_*}\eta_c\) is \(d_j\,o(b)\)
on the two type primes together with one threshold.
\end{proposition}

\begin{proof}
\proofstep{Use the label projector before applying the common tame pushout bound.}
For each good \(\pi\), form a point-dependent central
type operator on the cut space. It projects to the
series labels of \(\pi\) at both type primes, multiplied
by the row group order \(h\). The scalar action at \(\pi\)
uses the common isobaric partition for its Hecke
character. The support \(Y_\pi\) of this operator
contains \(\pi\). The agreement of the two series labels on the support gives
\[
 \pi\in Y_\pi,\qquad Y_\pi\cap Z_i\subseteq Z_i^+.
\]
This comparison
uses the finite series projectors, their realization on
the point spaces, and the span of those spaces in
each generic eigenspace. The operator acts integrally
and kills the eigenspaces outside \(Y_\pi\).

On \(Y_\pi\) the squarefree labels exclude the two
prescribed character values. The local--global
comparison supplies zero monodromy and the same
tame inertia type throughout this family. The inertia
polynomial of the full representation includes the
rank-\(t\) scalar character block. The corresponding
tame characters are inverted at the conjugate place.
Apply the type-prime period lifting comparison to the
twisted family on the positive side and to its Tate
dual on the negative side. The dual sends a jump \(a\) to \(-a-1\). Thus in both applications the jump parameters are
\[
 (\text{lower parameter},\text{upper parameter})=(\eta+t+1,s_T+S_T).
\]
The base field is \(\mathbb Q_\ell\) at every such use.
For \(N=R+1\), put
\[
 A_0=\lfloor N^2/4\rfloor.
\]

The integral period lift has loss
\[
 e_*=C_T\{N A_0+(\eta+t+1)(N+A_0)\}.
\]
The local pushout comparison sends this to normalized
length at most \(2u e_*\) modulo the finite
cohomology condition in rank \(u\). For the rank-\(r\)
constituent use \(u=r\), and for each of the \(t\)
character constituents use \(u=1\).
The same exponent applies throughout the locus \(Y_\pi\).
\label{cor:special-prime-lengths-5}

The global marked congruence comparison on \(Y_\pi\)
now gives an element of the Hecke order supported
on \(Z_i\cup(X_i\setminus Y_\pi)\). Its value at
\(\pi\) has valuation bounded by \(4fR e_*\)
plus the signature, strict cohomology, class-number,
mark and tame allowances. The local term is
\(2f(2re_*+2te_*)\) since all completions have
degree one. The marked determinant is evaluated at the very same normalized specialization. Over all places above \(\ell\) the local rank and character terms sum with degree factor
\[
 \sum_{w\mid\ell}[E_w:\mathbb Q_\ell]=[E:\mathbb Q]=2f.
\]

In this application the row
quotients are the local pushout images. With the
finite condition imposed, the three global kernels
are the signature, unramified class-number and
strict Tate-twist kernels. For the character block,
the coefficient character after twisting is
\(\chi_\ell^{\,1-c}\) at index \(c\).
The divisible finite image is unramified at \(c=1\).
It vanishes at \(c>1\) by the finite cohomology
formula for negative Tate twists. These conditions
give the class-number and strict kernels after
the tame restrictions away from \(\ell\).
The marked and tame costs
are \(4R v_\delta\) and
\[
 t\sum_{v\in\mathcal Q_i}v_\ell(\zeta_v-1).
\]

Multiplying the Hecke element by the integral
type operator gives support on \(Z_i^+\).
Apply the lattice filtration comparison with this
product for the graded piece at \(\pi\).
The other graded pieces may use different type
operators.

Since \(A_0\le R^2\) and \(N\le2R\), we have
\[
 e_*\le8C_TR^2(R+\eta+t).
\]

The valuation of the row group order at \(g=g_i\)
is also explicit. Its factor from \(g\) is
\(f\binom r2\) by the order formula for
\(\mathrm{GL}_r(\mathbb F_g)\). The factor from
\(g'=g_{1-i}\) is at most \(fr(r+1)\), because
\(g'^j-1<g^{2j}\) and hence
\[
 v_g(g'^j-1)\le2j.
\]
Each of the factors has
exponent \(f\), one place for each conjugate pair.
Their total valuation is therefore at most
\[
 f\binom r2+f\sum_{j=1}^r2j=\tfrac12 f(3r^2+r)\le2fr^2 .
\] Thus the whole
localization length \(\eta_i\) is bounded by the
multiplicity-weighted sum, over good \(\pi\), of
\[
 4fR\,C_T'R^2(R+\eta+t)+2fr^2
   +s_\pi+V_\ell+v_\ell(h_E)+4R\,v_\delta
   +t\sum_{v\in\mathcal Q_i}v_\ell(\zeta_v-1).
 \tag{T}\label{eq:local-prime-t}
\]
The constants in this inequality can be
taken to be \(C_T=40\) and \(C_T'=8C_T\).
Here \(v_\delta\) measures the cyclotomic mark of the
localized order. Its depth is one at either type
prime, so the marked determinant gives
\[
 v_\delta\le(t-1)+v_\ell((t-1)!)\le2t.
\]

The two killing allowances give
\(s_\pi\le\ell^{-2}\) and
\[
 t\sum_{v\in\mathcal Q_i}v_\ell(\zeta_v-1)\le\ell^{-2}.
\]

The cyclotomic invariant term vanishes because the
host is unramified at the odd prime \(\ell\) and \(tr\) is odd.
For a lift of \(-1\) the scalar difference is
\[
 (-1)^{tr}-1=-2\in\mathbb Z_\ell^\times .
\] For the signature
allowance, take the character of the chosen signature of
\(\pi\) in Section~\ref{subsec:signature-selmer} and its
specialization at the auxiliary character of this coefficient
place. The specialization is chosen simultaneously for the
good signatures. The lattice comparison identifies the
rank-\(r\) piece of the same marked representation with the
specialized signature lattice. The Selmer injection and the
normalized coefficient-extension comparison of the retained
twist give finiteness and transfer the length allowance.
The full-image generator of auxiliary inertia gives the
tame allowance on this same localization. The strict
cohomology lengths for \(1\le e<t\) are the separate
finite lengths in the definition of \(V_\ell\).

The uniform bounds used in the estimate are
\[
 R\le2t,\qquad t\le\eta_c,\qquad n,t\le x^{C_1}.
\]
Hence
the first, second and sixth terms of \eqref{eq:local-prime-t} have sum at most
\[
 (256C_T'+20)nt^3\eta_c
 \le (256C_T'+20)x^{4C_1}\eta_c .
\]
Put \(C_* = 2000(C_T'+1)+5C_1\). For \(x\ge1\) and
\(C_1\ge1\) the last bound is at most
\(C_*x^{C_*}\eta_c\). Thus the uniform bound per
copy of a good character at \(\ell\) is
\[
 M_\ell=C_*x^{C_*}\eta_c+
            2\ell^{-2}+V_\ell+v_\ell(h_E).
\]
\proofstep{Weight the lengths by the actual good multiplicities and the local degree identity.}
The local inequality applies to a single good character. In the numerator each occurs with its actual eigenspace multiplicity; the bound for their sum uses the dimension of the same selected row.

Use the local comparison including the empty case
\(\eta_i=0\) when \(Z_i^+=\varnothing\).
The ties defining the twisted row identify \(\eta_{j,w}\)
with \(\sum_{i\ {\rm over}\ (j,w)}\eta_i\). Apply
\eqref{eq:retained-multiplicity} to this sum:
\[
 \eta_{j,w}\le d_j M_\ell,\qquad
 \sum_{\substack{w\mid\ell\\w\ {\rm twisted}}}
        e_w\omega_w\eta_{j,w}\le d_j(\log\ell)M_\ell .
\]
The second inequality uses the degree sum over all places
above \(\ell\). These bounds use the actual multiplicities
and require no uniform upper bound for the number of
localizations.

In particular the contribution of the first, second and
sixth terms to the numerator, summed over both rational
type primes, is at most
\[
 d_j C_*x^{C_*}\eta_c(\log g_0+\log g_1)=d_j\,o(b),
\]
since \(\log g_i=O(\log b)\), \(\log\eta_c\le x^{\epsilon/4}\)
and \(b=\exp(x^\epsilon)\). The \(o(b)\) bound per unit of rank
is uniform on the chosen windows for fixed \(\epsilon\). The division
comparison uses \(d_j=D_0\) before dividing by \(D_0\).
The group-order factor in \((T)\) is already included at
these two primes. The separate cost of returning to the
uncut length and the common budgets for strict cohomology
and class number enter the full numerator comparison.
\label{cor:special-prime-lengths-6}
\end{proof}

\subsubsection*{The two ramified auxiliary primes}\label{subsec:ramified-auxiliary-row}

Fix \(\ell\in\{q,d\}\) and let \(\ell'\) denote the other prime.
The three factors left after omitting \(\sqrt{-\ell}\) from
\eqref{eq:host-character-compositum} have a common cyclotomic modulus
prime to \(\ell\). In fact, once the lower endpoint in
Proposition~\ref{prop:quadratic-prime-pair} is large, both auxiliary
primes exceed the fixed conductor \(f_\psi\). For each \(w\mid\ell\) put
\[
 M_\ell=\operatorname{lcm}(2^{e+2},f_\psi,\ell'),\qquad
 B_\ell=K_\psi C_e(\sqrt{-\ell'}),\qquad F_w=(B_\ell)_{w|_{B_\ell}}.
\]
The completion embeds in \(\mathbb Q_\ell(\mu_{M_\ell})\). Roots of
unity of order prime to the residue characteristic generate an
unramified extension of a complete discretely valued field
\cite[Corollary~10.16]{sutherland-local-fields}. In this application
the base is \(\mathbb Q_\ell\) and \(\ell\nmid M_\ell\). Hence
\begin{equation}\label{eq:host-auxiliary-completion}
 F_w/\mathbb Q_\ell\text{ is unramified},\qquad
 E_w=F_w(\sqrt{-\ell}),\qquad e(E_w/\mathbb Q_\ell)=2 .
\end{equation}

Indeed \(X^2+\ell\) is Eisenstein over \(F_w\), since
\(v_{F_w}(\ell)=1\). Since
\(\ell\equiv3\bmod4\), the unique quadratic subfield of
\(\mathbb Q_\ell(\mu_\ell)\) is \(\mathbb Q_\ell(\sqrt{-\ell})\)
by the quadratic Gauss sum. Writing \(\chi_{\rm cyc}\) for the cyclotomic character, we obtain
\[
 [F_w(\mu_\ell):E_w]=\frac{\ell-1}{2},
 \qquad
 \chi_{\rm cyc}(I_{E_w})\bmod\ell\ \supseteq\
           (\mathbb F_\ell^\times)^2 .
\]
The size choice gives \((\ell-1)/2>tr\). For every nonzero tested exponent \(a\),
\[
 0<|a|\le tr<(\ell-1)/2.
\]
The character with exponent \(a\) cannot be trivial on the square subgroup, which has order \((\ell-1)/2\). Thus \(\nu_\ell=0\).

The source describes the CM type by the sign of
\(\sqrt{-d}\). At \(q\), the congruence on the auxiliary discriminant gives
\[
 -d\equiv1\pmod q,\qquad -d\in\mathbb Q_q^{\times2}.
\]
Thus its decomposition group preserves the sign of \(\sqrt{-d}\). At \(d\) the nonidentity inertia element reverses it.
In either case the local jump list is taken at each
embedding, allowing its sign to depend on the
embedding of \(E_w\). These primes are distinct from
the two type primes, which split completely in the host.
The integral global congruence comparison away from
type primes applies. The declaration used in this ramified comparison is
\[
 \mathcal R_\ell=\{w:w\mid\ell\},\qquad \mathcal C_\ell=\varnothing.
\]

This choice requires a bound on each row at \(d\)
with mixed embedding signs. A count of rows obtained
by taking one sign for the whole completion would
lose the role of the inertia involution at \(d\).
\label{cor:special-prime-lengths-7}

At each embedding the original jumps are
\(0,\ldots,t-1\), \(m\) central jumps and \(k-1\) jumps
in a signed interval of width \(S_c\), separated from
the center by more than \(\eta_c+t+1\).
Put \(J_\tau\) for the original outer integer jump set at
embedding \(\tau\). The Tate operation on the outer jumps is
\[
 J_\tau\longmapsto -1-J_\tau=\{-1-j:j\in J_\tau\},
\]
with
\[
 \begin{aligned}
 \max(-J_\tau-1)-\min(-J_\tau-1)
    &=\max J_\tau-\min J_\tau\le\lfloor S_c\rfloor,\\
 \min_{a\in J_\tau}|-a-1|
    &\ge\min_{a\in J_\tau}|a|-1>\lfloor\eta_c\rfloor+t .
 \end{aligned}
\]
Both the original jump list and its Tate-dual image have
absolute bound \(\lfloor2b\rfloor\). Put \(I=[-\lfloor\eta_c\rfloor-t,\lfloor\eta_c\rfloor+t]\). For each embedding \(\tau\), removal of the cyclotomic summand leaves exactly the simple central jumps of the rank-\(r\) dual, so
\[
 \#\{a\in\mathrm{Jumps}_\tau(T_2):a\in I\}=m.
\]

For the row source, a subrepresentation of a sum of
the full duals is crystalline and has its jumps in
their union. This uses strictness of the crystalline
filtration and the common jump set at the embedding.
These are the family and target hypotheses of the
integral row comparison.

The local result used here is the crystalline
row-denominator theorem with ramification two.
For a surjection of stable lattices as above with
rank-\(r\) target, let the source jumps at each
embedding lie in the central interval of radius
\(\eta_0\) together with a signed interval of width
\(S\), lying at distance greater than \(\eta_0\)
from zero. Let all jumps have absolute value at
most \(h_0\). The target has at most \(m\) jumps,
counted with multiplicity, in the central interval.
Then the normalized finite-cohomology colength
for this surjection is at most
\[
 160\bigl((D_w+1)r(2h_0/\ell+2S+
             2\log_\ell h_0+1+\eta_0)
                   +3D_w m(h_0+1)\bigr).
\]
We derive this estimate from the integral period comparisons in the
next paragraphs. The coefficient valuation throughout has
\(v(\ell)=1\). \label{cor:special-prime-lengths-8}

\paragraph{The integral maps at one cyclotomic level.}
\label{par:row-metric}
\hypertarget{source-row-metric}{}
\begin{paramdefinition}[Local field, signed intervals and integral map]\label{hyp:first-level-row}
Write \(p=\ell\), let \(K/\mathbb Q_p\) be the local field, and let
\(K_0\) be its maximal unramified subfield. Its residue degree,
ramification index and absolute degree are
\[
 n=[K_0:\mathbb Q_p],\qquad e=[K:K_0],\qquad D=[K:\mathbb Q_p]=ne.
\] In the application \(e=2\).
The calculation also applies at \(e=1\).
A first-level row has odd residue characteristic,
\[
 K\subset K_0(\mu_p),
\]
and a coefficient
DVR \(O\) containing the images of the rings of integers
under all embeddings of \(K_0(\mu_p)\). All the following lengths are
divided by \({\rm length}_O(O/pO)\). Let \(F:T_1\twoheadrightarrow T_2\)
be the given integral crystalline surjection, with \(r={\rm rank}_O T_2\).
Write \(h,S,\eta\) for the integral height, width and central radius.
If the sign at \(\tau\) is \(d_\tau\in\{+1,-1\}\), the corresponding outer jumps \(a\) satisfy
\[
 b_\tau\le d_\tau a\le b_\tau+S.
\]

The signed jump data include a central interval \([-\eta,\eta]\) and a bound \(m\) on the number of central target jumps at each embedding. Call these data a first-level row when \(\eta<b_\tau\) at every embedding.
\end{paramdefinition}

For the two ramified auxiliary primes \(p=q,d\), \eqref{eq:host-auxiliary-completion} gives \(K=K_0(\sqrt{-p})\) with \(p\equiv3\pmod4\). The quadratic Gauss sum gives
\[
 K=K_0(\sqrt{-p})\subseteq K_0(\mu_p),\qquad [K:K_0]=2 .
\]
At an unramified place \(K=K_0\). A finite coefficient extension contains all the indicated embedding images. All actual jumps lie in \([-h,h]\); clipping an empty outer interval or the central interval outside this range therefore preserves support and separation. The normalized endpoints are
\[
 \eta' = \min(\eta,h),\qquad b'_\tau=\min(b_\tau,h+1),
 \qquad \eta'<b'_\tau\le h+1 .
\]
Use the original symbols for these normalized values. Then \(H=h+S+1\) satisfies \(b_\tau+S\le H\). Strict exactness of the crystalline filtration puts the target jumps in the same support.

We describe the integral comparison used here. At full level one
there are \(n\) components indexed cyclically by the embeddings of
\(K_0\). For the tame group \(G\), one has
\[
 |G|=s=(p-1)/e,\qquad O[[\pi]]^G=O[[X]],\qquad
 X\in\pi^s O[[\pi]]^\times,\qquad
 \frac{1}{s}\sum_{g\in G}g\in O[G].
\]
The last element is the integral projection to its invariants.

Let \(E=E_c(X)\)
be the degree-\(e\) polynomial associated to the first cyclotomic
factor on component \(c\), up to a unit. Its roots are the images of
\(X\) at the first cyclotomic points. They belong to the maximal
ideal of \(O\). We use a generator \(\gamma\) of the principal
cyclotomic direction, \(u=\chi_{\rm cyc}(\gamma)\), \(v(u-1)=1\).
On the invariant ring
\[
 \gamma X=u^s X\bmod X^2
\]
and
\(\varphi X\) is a unit times \(X E^s\), using the target component
of Frobenius.

\(N_c(T)\) is the tame-invariant Wach lattice and \(P_c(T)\)
is the invariant Frobenius span from the preceding component. Both
are free over \(A\), with
\[
             E^H P_c\subset N_c\subset E^{-H}P_c .
\]
The action of \(\gamma\) on \(N_c/XN_c\) is trivial. After inverting
\(p\) this quotient gives the corresponding crystalline constants.
Frobenius identifies the constants from \(c-1\) with those used
to evaluate \(P_c\).

\begin{lemma}[Filtration test on bounded numerators]
On each labelled first-level component with logarithmic parameter
\(t\), write \(a\) for the numerator of the element \(E^{-H}a\in E^{-H}P_c\). The test is
\[
 E^{-H}a\in N_c\quad\Longleftrightarrow\quad [t^i]a\in {\rm Fil}^{H-i}_\tau\quad\text{for every first-level root and }i\ge0.
\]
\end{lemma}
\begin{proof}
At a primitive root use a logarithmic parameter
\(t\) on which \(\gamma t=ut\). On bounded-pole Frobenius coordinates,
membership in \(N_c\) is tested by requiring the coefficient of
\(t^i\) to lie in \({\rm Fil}^{-i}_\tau\) at every point of the
component. In numerator coordinates the same condition for a vector \(z=E^{-H}a\) becomes
\[
 [t^i]a\in{\rm Fil}^{H-i}_\tau\qquad(\text{each degree }i).
\]
Indeed \(E\) has a simple zero, and multiplication by a power-series unit preserves the coefficient tests for a decreasing filtration. The evaluation identifies the rational constant
space with the labeled de Rham space using the crystalline comparison;
lifts of the same constant give the same constant value, and the
kernel on the preceding integral lattice is \(XN_{c-1}\).
One uses each embedding over \(c\), including both at ramification
two.

In an adapted de Rham basis of jumps \(a_j\), the completed
filtered lattice is generated by \(t^{-a_j}d_j\), which gives
the coefficient test. Its integral use on already bounded
numerators follows by division by the distinguished factors of
\(E\): if a remainder vanishes at the full set of roots with the
specified multiplicities over the coefficient field, the integral
monic remainder is zero. The height inclusions supply the initial denominator.
\label{cor:row-metric-1}
\end{proof}

We use the integral Wach comparison,
étale exactness, and the finite Herr class criterion for these
lattices. The finite criterion sends a tame-invariant restricted
cocycle pair \((w,z)\), with \(w\in E^{-H}P_c\) and \(z\) sufficiently
deep in \(N\) as below, to a finite class, naturally for \(F\).
The extension comparison sends the connecting pair of a lift of the
unit in a crystalline extension to its integral extension class
up to the fixed unit of the trivial frame.
\label{cor:row-metric-2}

Put
\[
 a=\lceil 2H/s\rceil,\quad M=\lceil4H/s\rceil,\quad
 A_i=i+v_p(i!),\quad \rho=A_{2a-1},\quad \sigma=A_{M-1}.
\]
\begin{proposition}[Integral comparison at the first cyclotomic level]
\label{prop:one-level-integral-row-maps}
For the local field, signed heights and maps in Definition~\ref{hyp:first-level-row},
the map on tame-invariant Wach lattices satisfies
\[
 p^\rho N_c(T_2)\subseteq F(N_c(T_1)).
\]
A first coordinate invariant modulo \(XN_c\) can be multiplied
by \(p^\sigma\) and corrected to become invariant modulo
\(X^MN_c\). It then has a
unique second coordinate in \(X^MN\). The two filtered coordinate
maps \(\alpha,\beta\) have composite \(p^{2\sigma}\) modulo the
specified error modules and preserve the tests modulo \(X^j\)
for \(1\le j\le M\). The same exponents work on the source,
target and unit extension.
\end{proposition}

\begin{proof}
\proofstep{From the tame average to a constant cokernel annihilator.}
First \(p^\rho N_c(T_2)\subset F N_c(T_1)\). For the initial parameter
bound use the ordinary Wach cokernel computation of
Section~\ref{subsec:ramified-lift} on the full cyclotomic field at level
one. Under the étale and Wach identification after restriction to
\(K_0(\mu_p)\), the underlying full-level map is the original map
with the tame action forgotten. The two height bounds of size \(H\)
give \(\pi^{2H}\) in its cokernel annihilator by extracting successive
Frobenius coordinates of a polynomial in \(1+\pi\).

If \(y\) belongs
to the tame-invariant target, \(X^a y\) therefore has a preimage in
the full-level component, since \(sa\ge 2H\). If \(Fz=X^a y\) and \(T_{\rm tame}\) is the tame group, then
\[
 z_0=|T_{\rm tame}|^{-1}\sum_{g\in T_{\rm tame}}gz,\qquad
 Fz_0=|T_{\rm tame}|^{-1}\sum_{g\in T_{\rm tame}}g(X^ay)=X^ay.
\]
The group order is a unit, so this gives an integral invariant preimage. This proves
\(X^a N_c(T_2)\subset F N_c(T_1)\).

On \(X^iN_c/X^{i+1}N_c\) the operator \(\gamma\) is the scalar
\(u^{si}\). This is true on the corresponding filtration of the
image as well, by equivariance. On \(X^ay\) for \(y\in N_c(T_2)\),
apply the two polynomials with roots \(u^{si}\) for \(0\le i<a\)
and \(a\le i<2a\). Both results lie in \(X^aFN_c(T_1)\), by the
image filtration for the first and the target filtration for the
second.

Put \(r_i=u^{si}\) and
\[
 L_i(Z)=\prod_{j\ne i}(Z-r_j)/(r_i-r_j)
\]
on the \(2a\) nodes.
The sum of \(L_i\) for \(i<a\) vanishes on the \(a\) upper
nodes, and that for \(i\ge a\) vanishes on the lower nodes.
Also \(\sum_i L_i=1\). The bound below gives \(p^\rho L_i\in O[Z]\),
so monic division by the products of the opposite roots
gives an identity
\[
         p^\rho=A(Z)\prod_{i<a}(Z-r_i)
                    +B(Z)\prod_{i\ge a}(Z-r_i),
                    \qquad A,B\in O[Z].
\]
Applying it at \(\gamma\) proves the claimed inclusion.

For distinct nodes \(v_p(r_i-r_j)=1+v_p(i-j)\), so the
denominator of the \(i\)-th interpolation polynomial costs
\[
 \sum_{j\ne i}v_p(r_i-r_j)
 =2a-1+v_p\bigl(i!(2a-1-i)!\bigr)
 \le 2a-1+v_p\bigl((2a-1)!\bigr)\le\rho .
\]
The first inequality uses the integrality of the binomial coefficient.
Cancel \(X^a\) to obtain the assertion. The same computation applies
to the surjection from an extension by \(T_2\) to the unit object,
whose extra jump is zero.

\proofstep{Raising the precision and completing the Herr pair.}
A second use of these scalars lifts invariance. If
\(w\in E^{-H}P_c\) satisfies \((\gamma-1)w\in XN_c\), one can
correct \(p^\sigma w\) by an element of \(XN_c\) so that its
difference under \(\gamma\) belongs to \(X^MN_c\). If \(R_i=(\gamma-1)w_i\in X^iN_c\), set
\(w_{i+1}=(u^{si}-1)w_i-R_i\). The scalar action on the \(i\)-th
graded piece gives the next precision:
\[
 (\gamma-1)w_{i+1}=(u^{si}-\gamma)R_i\in X^{i+1}N_c .
\]
The product over \(1\le i<M\) is a unit times \(p^\sigma\).
Also an element of \(XN_c\) invariant to order \(M\) already lies in
\(X^MN_c\), by its first nonzero graded coefficient and
\(u^{si}\ne1\).

On \(X^MN\) the operator \(1-\varphi\) is invertible
by its geometric series. Indeed \(sM-H\ge1\) leaves a factor \(E\)
after applying Frobenius and clearing its height denominator, so
iteration goes to zero in the complete \((p,\pi)\)-adic topology.
Consequently a first Herr coordinate invariant to order \(M\)
has a unique second coordinate in \(X^MN\) solving
\[
 (\varphi-1)z=(\gamma-1)w.
\]
Changing its first coordinate by
\(X^MN\) changes the pair by a coboundary, by the same inverse.

\proofstep{Transporting the horizontal bases by Frobenius.}
For later use we describe the coordinate maps together with their
precision. Let \(L=N_{c-1}/XN_{c-1}\). In an integral basis of
\(N_{c-1}\) lifting a basis of \(L\), construct a matrix \(U(X)\)
with \(U(0)=1\) so the corrected basis is \(\gamma\)-invariant
modulo \(X^M\) over the fraction field.

If the residual first occurs
in degree \(i\), the coefficient equation divides that residual by
\(1-u^{si}\). Thus the degree \(i\) coefficients of \(U\) have
denominator at most \(p^{A_i}\) for \(i<M\). The inverse has the
same bound, since
\[
 \sum A_{i_j}\le A_{\sum i_j}
\]
by factorial
divisibility. Frobenius transport of \(p^\sigma U\) and
\(p^\sigma U^{-1}\) furnishes two integral maps between the model
\(E^{-H} A\otimes_O L\) and \(E^{-H}P_c\), modulo
\(X^M E^{3H}\) times the respective undivided modules. Their composites are \(p^{2\sigma}\) in these quotients because
\[
 \frac{\varphi(X)^M}{E^H}\in X^M E^{sM-H}A
                         \subseteq X^M E^{3H}A
 \qquad(sM\ge4H).
\]
They commute in the quotients with \(\gamma\) and lattice maps.
For the latter property the map in horizontal bases agrees through
degree \(M-1\) with its constant matrix, by equivariance and
\(u^{si}\ne1\). Denote the forward and backward maps by \(\alpha,\beta\).

\proofstep{The filtration tests at all first-level roots.}
In the first-level expansion the corrected basis specializes to its
constant through orders below \(sM\) before division by \(E^H\).
Indeed the pulled-back invariant equation forces all intermediate
positive coefficients to vanish, since the degree \(i\) cyclotomic
eigenvalue is \(u^i\ne1\). The precision and jump-height bounds give
\[
 i\ge sM\ \Longrightarrow\ i\ge4H\ \Longrightarrow\ i\ge2H.
\]
All first-level filtration tests for numerator coefficients in those degrees are automatic. Hence both maps carry the filtered model lattice,
denoted \(N^{\rm mod}\), into \(N_c\) and back, respectively, modulo
the indicated error modules. They preserve the tests for membership
in \(X^jN\) for \(1\le j\le M\) in the forward or backward direction
from a member, by linearity in the invariant coefficient ring and
the inclusion of the error in \(X^MN\). In particular an invariant
modulo \(XN\) has an invariant partner. These constructions apply
with the same \(H,M,\rho,\sigma\) to the source, target and unit
extension.
\label{cor:row-metric-3}
\end{proof}

\paragraph{The number of exceptional coordinates.}
\label{par:row-restricted}
At component \(c\) of the target, use the lattice
\(L=N_{c-1}/XN_{c-1}\) of Section~\ref{par:row-metric}.
Identify \(L\otimes_O{\rm Frac}(O)\)
with the crystalline constant space and pull every labeled filtration
back to it. At a plus embedding, \(C_\tau={\rm Fil}^{b_\tau}_\tau\)
has codimension at most \(m\). At a minus embedding,
\(B_\tau={\rm Fil}^{1-b_\tau}_\tau\) has dimension at most \(m\).
These inequalities follow by counting the central jumps, including
their multiplicities, on the two sides of the respective cuts.
The filtered spaces for distinct roots belonging to the same embedding
give the same test on the crystalline constants.

\begin{proposition}[A bounded set of exceptional coordinates]\label{prop:row-exceptional-coordinates}
Assume \(\operatorname{codim} C_\tau\le m\) at the plus
embeddings and \(\dim B_\tau\le m\) at the minus embeddings,
with at most \(e\) embeddings in total. There is a direct summand \(B\) with complement \(L_{\rm main}\) such that
\[
 L=L_{\rm main}\oplus B,\quad L_{\rm main}\subseteq\bigcap_\tau C_\tau,\quad B_\tau\subseteq B\otimes K\quad(\tau\text{ minus}),\quad \operatorname{rk}_O B\le em,
\]
where \(K\) is the fraction field and the intersection is over the plus embeddings.
The image \(I_B\) of the invariant residual coordinates defined
below is free of rank at most \((e+1)\operatorname{rk}_O B\).
\end{proposition}

\begin{proof}
\proofstep{Splitting off the exceptional directions.}
There is an integral splitting \(L=L_{\rm main}\oplus L_{\rm err}\)
such that \(L_{\rm main}\) lies in every \(C_\tau\) and
\(L_{\rm err}\otimes{\rm Frac}(O)\) contains every \(B_\tau\), with
\({\rm rank}\,L_{\rm err}\le em\). Let \(K=\operatorname{Frac}(O)\), \(C=\bigcap_{\tau,+}C_\tau\),
and write \(L=(C\cap L)\oplus D\) using saturation. Define
\[
  L_{\rm err}=L\cap\Big(D_K+\sum_{\tau,-}B_\tau\Big),
  \qquad
  \operatorname{rk}L_{\rm err}\le \operatorname{codim} C
                          +\sum_{\tau,-}\dim B_\tau\le em .
\]
Choose \(L_{\rm main}\) complementary to \(L_{\rm err}\cap C\)
in \(C\cap L\). Then \(L=L_{\rm main}\oplus L_{\rm err}\)
since \(D\subseteq L_{\rm err}\). All splittings used here exist over the DVR because
a saturated submodule of a finite free module has free quotient.
In particular the projection to \(L_{\rm main}\) has integral entries
and kills all the \(B_\tau\).

\proofstep{The polynomial that separates the main coordinates.}
Consider a model coordinate \(z\in E^{-H}A\otimes L\) with
\[
 (\gamma-1)z\in XN^{\rm mod},
\]
for example \(\beta(w)\) with \((\gamma-1)w\in XN_c\). At a primitive root the coefficient of
\(t^i\) in \(z\), for \(i\ne0\), belongs to \({\rm Fil}^{-i}_\tau\).
Indeed the coefficient of its difference is multiplied by \(u^i-1\),
nonzero in the fraction field, and multiplication by \(X\), a unit
at this root, preserves the filtration inequalities. On projecting
to \(L_{\rm main}\) it follows that the coefficients vanish in degrees
\(i<-b_\tau-S\) at plus embeddings, and in nonzero degrees
\(i<b_\tau\) at minus embeddings. In the latter case the coefficient
is in \(B_\tau\), which the projection kills.

Let \(\Omega\) be the union over these embeddings of the integer
intervals \([-b_\tau-S,-b_\tau)\) at plus and
\([b_\tau,b_\tau+S)\) at minus. The exceptional degrees determine the polynomial
\[
 P(Z)=\prod_{i\in\Omega}(Z-u^i).
\]
The operator \(P(\gamma)\) preserves \(E^{-H}A\) since \(\gamma E/E\)
is a unit. In the Laurent parameter at that root, write the projection
to the main summand as \({\rm pr}_{\rm main}z=\sum_i z_i t^i\).
Since this projection commutes with \(\gamma\),
\[
 [t^i]{\rm pr}_{\rm main}P(\gamma)z=P(u^i)z_i,\qquad
 P(u^i)=0\quad(i\in\Omega).
\]
Below the indicated intervals the projected coefficients already vanish,
with the possible exception of the constant at a minus embedding.
The polynomial removes the degrees inside the intervals. Write the main part of \(P(\gamma)z\) in the integral
basis obtained above as \(\sum f_j\otimes l_j\).

At plus points \(f_j\) has order at least \(-b_\tau\); at minus
points
\[
 f_j=x_{\tau,j}+O(t^{b_\tau+S})
\]
for a constant
\(x_{\tau,j}\) in the coefficient field (with the value at the
chosen root understood). Also \((\gamma-1)f_j\in X E^{-H}A\).
The latter inclusion follows by commuting projection and the polynomial
with \(\gamma\) and using the stability of \(XE^{-H}A\).
Each simple tensor is invariant modulo \(XN^{\rm mod}\).
For its difference divided by \(X\), the required filtration at
plus points follows from \(l_j\in{\rm Fil}^{b_\tau}\) and the lower
order just proved. At minus points the constant cancels and all
of \(L\) is in \({\rm Fil}^{-b_\tau-S}\).
Thus the residual part is also invariant modulo \(XN^{\rm mod}\).

Finally
\[
 P(\gamma)z\equiv P(1)z\bmod XN^{\rm mod},
\]
by stability
and \((Z-1)\mid(P(Z)-P(1))\).
The scalar cost to replace \(P(1)\) by a power of \(p\) is at most
\[
             \kappa=e\{S+v_p(S!)+\lfloor\log_p H\rfloor\}.
\]
Each interval has \(S\) nonzero integers of magnitude at most \(H\).
For one interval \(I\), the principal-unit formula therefore gives
\[
 \sum_{i\in I}v(1-u^i)
 = |I|+\sum_{a=1}^{\lfloor\log_p H\rfloor}
                  \#\{i\in I:p^a\mid i\}
 \le S+v_p(S!)+\lfloor\log_p H\rfloor .
\]
There are at most \(e\) intervals; overlaps can only reduce
the sum over their union. An empty union has cost zero.

\proofstep{Counting the residual directions.}
We need the number of integral residual classes, including the
coefficient from the ring variable. For \(B=L_{\rm err}\), consider
the image \(I_B\) in \((E^{-H}A\otimes L)/XN^{\rm mod}\) of those
\(z\in E^{-H}A\otimes B\) with \((\gamma-1)z\in XN^{\rm mod}\).
The two assertions to prove are torsion-freeness and the bound
\[
 \operatorname{rank}_O I_B\le(e+1)\operatorname{rank}_O B.
\]
Over the coefficient DVR they imply the asserted freeness once finiteness is known.
The ambient quotient is finite over \(O\): \(N^{\rm mod}\) contains
\(E^HA\otimes L\) by the jump bounds, and \(XE^{2H}\) is a unit times
a distinguished polynomial.

The ambient quotient is also torsion-free. If \(a z\in
XN^{\rm mod}\) for \(0\ne a\in O\), the constant coefficient of the
integral numerator of \(a z\) at \(X=0\) vanishes. Thus \(z/X\) still
has an integral numerator. At every primitive root the coefficient test for \(az/X\) lies in a
filtration space over \(K=\operatorname{Frac}(O)\). For each degree \(i\),
\[
  a[t^i](z/X)\in {\rm Fil}^{-i}_\tau
      \quad\Longrightarrow\quad [t^i](z/X)\in {\rm Fil}^{-i}_\tau .
\]
Together with integrality at zero this gives \(z/X\in N^{\rm mod}\).
The ambient quotient is consequently finite torsion-free over the DVR.
Its submodule \(I_B\) is free.

For the rank, use the finite filtration starting with \(E^{-H}A\otimes B\).
Multiply successively by \(X-r_j\), \(2H\) times for each of the
\(e\) roots \(r_j\) of \(E\), and finally by \(X\). The last term
is contained in \(XN^{\rm mod}\). Every grade has rank at most
\({\rm rank}\,B\), by division with remainder at its root, including
the last root zero. The filtration is \(\gamma\)-stable. At \(r_j\)
the local parameter is a unit times the cyclotomic logarithm. To see
this, write
\[
 X=\pi^s v(\pi)
\]
with \(v\) a unit.

At
\(\pi=\zeta-1\) its derivative is
\(\pi^{s-1}(sv(\pi)+\pi v'(\pi))\ne0\), since \(s\) is a \(p\)-adic
unit. The \(p-1\) simple roots in the ordinary parameter give \(e\)
simple roots in the invariant ring: the tame orbits have size \(s\)
and \(E(X(\pi))\) is a unit times the cyclotomic polynomial at
\(1+\pi\). The identical multiplicity calculation also shows that
different orbits give distinct roots. Alternatively \(E\) has
degree \(e\) and its derivative is nonzero at each image, and division
over each fibre of \(X-r_j\) in the ordinary disc has degree \(s\).

On the grade with preceding multiplicity \(a\) at \(r_j\), \(\gamma\)
acts by \(u^{a-H}\); factors coming from other roots are units there.
Since \(u\) has infinite order, the invariant-grade test at this root is
\[
 u^{a-H}=1\quad\Longleftrightarrow\quad a=H.
\]
The
grade at zero contributes at most one more copy of \(B\).
Over the fraction field filter the image in the quotient by
\(XN^{\rm mod}\). Its graded pieces are quotients of the original
grades, and a fixed vector has a fixed leading term. The action in the graded factors at \(r_j\) gives
\[
 \operatorname{rk}_O I_B
 \le \sum_{j=1}^{e}\sum_{\substack{0\le a<2H\\u^{a-H}=1}}
           \operatorname{rk}_O B+\operatorname{rk}_O B
 \le (e+1)\operatorname{rk}_O B\le e(e+1)m .
\]
Thus an \(O\)-basis of \(I_B\) furnishes a single choice of
representatives at this component for all target classes. Taking a different basis for each input
would give no bound on the number of exceptional generators.
\label{cor:row-coordinates-1}
\end{proof}

\paragraph{Lifting an outer coordinate through a quadratic component.}
\label{par:row-mixed}
\hypertarget{source-row-mixed}{}

\begin{lemma}[Lifting the main tensors]\label{lem:mixed-row-tensors}
For the main tensors on a one-sign component, or on a mixed
component with \(e=2\), put
\(y_c=\lceil\eta+1/2+\log_pH\rceil\).
For each such tensor, let \(y\in N_c(T_2)\) represent its image
under the forward comparison \(\alpha\) of Section~\ref{par:row-metric},
modulo \(XN_c(T_2)\). There exists \(v\in N_c(T_1)\) with
\[
 (\gamma-1)v\in XN_c(T_1),\qquad Fv\equiv p^{\rho+\max(2\sigma,y_c)}y\pmod{XN_c(T_2)}.
\]
\end{lemma}

\begin{proof}
Let \(f\otimes x\) be one of the simple tensors above. We require an
invariant in the source with the same image as a power of \(p\) times
its image under \(\alpha\) modulo \(XN_c(T_2)\). Two operations suffice.
First, if an invariant model tensor has an invariant source partner
under \(F\), apply \(\alpha\) and use naturality. Changing the partner by \(u\in X^M E^{3H}P_c\) preserves its cut class because
\[
 X^M E^{3H}P_c\subseteq XN_c.
\]

Second, if a model
tensor belongs to \(N^{\rm mod}_2\), its image under \(\alpha\) belongs to
\(N_c(T_2)\), hence has an integral lift after multiplication by
\(p^\rho\). Any such lift is invariant modulo \(XN_c(T_1)\).
Here a model tensor with a scalar coefficient and a preceding constant
can be lifted by multiplying the scalar by a lift of that constant,
provided the resulting tensor meets the invariance test. Actual
preceding constants have lifts after \(p^\rho\) by the same Wach
cokernel bound on the preceding component.

\proofstep{Components with one sign.}
At a component with only plus tests, the product of each
coefficient of \(f\) with \(x\) lies in the required
filtration, since \(f\) starts in degree \(-b_\tau\) or
later and \(x\in\operatorname{Fil}^{b_\tau}\). Thus
\[
 f\otimes x\in N^{\rm mod}_2.
\]

At a component with only minus tests, choose a preceding constant
\(x_1\) with
\[
 F x_1=p^\rho x,
\]
and take \(f\otimes x_1\).
The source conditions are
\[
  E^H f\in A,\qquad
  (\gamma-1)f\in XE^{-H}A,\qquad
  \operatorname{ord}_{\tau}\big((\gamma-1)f/X\big)
      \ge b_\tau+S .
\]
The last inequality suffices because the entire source constant
space is in \(\operatorname{Fil}^{-b_\tau-S}\) at each minus test.

\proofstep{The two roots of the quadratic component.}
For opposite signs assume \(e=2\). Write \(r_U,r_D\) for the
two parameter values at the plus and minus tests, and put
\[
 y_U=X-r_U,\ y_D=X-r_D.
\]
Division by \(X-r\) at any such
value is integral since \(r\) belongs to the maximal ideal and
\(O\) is complete. Over the fraction field the changes to \(y_U\) and \(y_D\) have invertible linear term at their respective roots. Each preserves the order of vanishing and transports the filtration test. In particular equal \(X\)-values impose identical scalar Taylor conditions.

If a plus test has the minus
root as well, the high vanishing of \((\gamma-1)f/X\) at the
minus root also suffices at that test, since its positive degree
is greater than \(\eta\). This gives the preceding-lift argument
when every plus test shares a minus root. Otherwise there are just
two distinct roots, the minus tests at \(r_D\) and the remaining
plus tests at \(r_U\), with the same labeled filtration at all
points over the corresponding embedding.

Write
\[
 b=b_U,\qquad B=b_D+S,\qquad
 E=a_2 y_U y_D,\qquad v(a_2)=0,\qquad d_*=r_D-r_U.
\]
Both roots have valuation \(s/(p-1)=1/2\). Also
\(v(d_*)=1/2\): at \(r_U\) differentiate the cyclotomic identity
\(Q(\pi)={\rm unit}\cdot E(X(\pi))\), where
\(Q=((1+\pi)^p-1)/\pi\). The derivative of \(Q\) at a primitive
point has valuation \(1-1/(p-1)\), and \(X'(\pi)\) has valuation
\((s-1)/(p-1)\). Hence
\[
 v(E'(r_U))=v(Q'(\pi))-v(X'(\pi))=1-\frac1{p-1}-\frac{s-1}{p-1}=\frac12.
\]
This accounts for the half power in the following allowance.
Let \(y_c=\lceil\eta+1/2+\log_p H\rceil\).

\proofstep{Nearby breaks.}
Suppose first \(b\le B+2\eta\). The numerator \(a=E^H f\) is
divisible by \(y_U^{H-b}\) in \(A\). At the down root the series
of \(f\) equals a constant \(x_0\) modulo degree \(B\).
Consequently the coefficients of \(a/y_U^{H-b}\) there, through
degree \(H+B-1\), agree with those of
\(x_0 a_2^H y_D^H(y_D+d_*)^b\). They are integral in \(y_D\)
by division in \(A\). In particular
\[
           x_0a_2^H \binom{b}{j} d_*^{\,b-j}\in O
                   \quad (0\le j<B,\ j\le b).
\]

If \(b<B\) take \(j=b\). Otherwise take \(j=B-1\); then
\(b-j\le2\eta+1\) and
\[
 v_p\binom{b}{j}\le\lfloor\log_p b\rfloor
\]
by the factorial formula, a difference of at most one in each
power-\(p\) count. In either case \(p^{y_c}x_0\in O\).
Subtract this constant from \(p^{y_c}f\). The remainder times
\(x\) belongs to \(N^{\rm mod}_2\). At the minus points its
order is at least \(B\); at the plus points its order is at
least \(-b\) and the target constant is in \({\rm Fil}^b\).
The constant term itself lifts by an arbitrary preceding lift,
since a constant is fixed by \(\gamma\). Applying the two operations
gives an invariant lift in the source Wach lattice at exponent \(\rho+y_c\).

\proofstep{Separated breaks.}
Suppose now \(B+2\eta<b\). A division by \(d_*^{b-B+1}\) could
cost an amount proportional to the outer height. Instead form the
auxiliary numerator
\[
 a_0=y_U^{H-b}y_D^{H+B}.
\]
The tensor
\(a_0 x/E^H\) belongs to \(N^{\rm mod}_2\) by its orders at both
roots. Apply \(\alpha\), lift its image in the target Wach lattice after \(p^\rho\)
in \(N_c(T_1)\), and apply \(\beta\).

Write the result with denominator
\(E^H\) as \(\sum_j a_j v_j/E^H\), \(a_j\in A\), where \(v_j\)
are integral preceding constants of the source. The sum in every
numerator coefficient of degree \(i\) at a root lies in
\({\rm Fil}^{H-i}_\tau\). On applying \(F\) its expansion in degrees
through \(H\) agrees with
\[
                  p^{\rho+2\sigma}a_0 x .
\]
Indeed the forward and backward composites give this scalar and the
coordinate modulus, after multiplying by \(E^H\), contains the
factor \(E^{4H}\) in every error term. Integral bases on the preceding
component give a representation in this common constant space for
both filtrations.

Expand each \(a_j\) at \(r_D\) through degrees \(<H-\eta\), with
coefficients in \(O\), and remove this polynomial in \(y_D\).
In degree \(i<H-\eta\) the corresponding sum against \(v_j\)
belongs to
\[
 {\rm Fil}^{H-i}_D=0,
\]
since the largest source
jump there is at most \(\eta\). These polynomial coefficients
therefore give zero sums against both \(v_j\) and \(Fv_j\) at the
plus root also, by the common crystalline constant identification.
Divide the remaining functions by \(y_D^{H-\eta}\), obtaining
integral \(a'_j\). At \(r_U\) this divisor is a unit over the
fraction field and the source coefficient condition
\({\rm Fil}^{H-i}_U\) persists. The new target expansion agrees
through degree \(H\) with
\[
       p^{\rho+2\sigma}y_U^{H-b}
                  (y_U+r_U-r_D)^{B+\eta}x .
\]

Take its coefficient of degree \(j_*=H-b+B+\eta\) in \(y_U\).
The assumption gives \(j_*<H-\eta\); changing from \(t\) to \(y_U\)
preserves the filtration test and precision, by triangular substitution.
This is the leading coefficient of the displayed polynomial. Hence
the same integral linear combination of the \(v_j\) maps to
\(p^{\rho+2\sigma}x\) and belongs to
\({\rm Fil}^{H-j_*}_U\subset{\rm Fil}^{\eta+1}_U
={\rm Fil}^b_U\). The equality uses the empty interval of jumps
strictly between \(\eta\) and \(b\). This combination supplies a
preceding lift compatible with the plus test. Keeping \(f\) gives
an invariant model lift at both signs, hence an invariant Wach-lattice
lift after \(\alpha\). Altogether a single exponent
\[
                       y_1=\rho+\max(2\sigma,y_c)
\]
serves every main tensor on every component, including the one-sign
cases.  \label{cor:row-mixed-1}
\end{proof}

\paragraph{From the Herr pairs to normalized colength.}
\label{par:row-colength}
On each component use the module \(I_B\) of exceptional coordinates
from Proposition~\ref{prop:row-exceptional-coordinates}. Choose lifts
of an \(O\)-basis of \(I_B\).
Apply \(\alpha\) to each lift, choosing a representative invariant
modulo \(XN_c\). By the order-raising argument in
Section~\ref{par:row-metric}, its \(p^\sigma\)-multiple can be
corrected modulo \(XN_c\) to an invariant to order \(M\).
Complete it by the unique deep second Herr coordinate. This gives
at most \(e(e+1)m\) fixed finite classes on this component.
Write \(H^1_f(K,T)\) for the inverse image of the rational
finite condition in integral cohomology, including torsion.
Write \(J\) for the module generated by all these classes on all
components, in the quotient
\[
              C_F=H^1_f(K,T_2)/F_*H^1_f(K,T_1).
\]
There are at most \(ne(e+1)m=D(e+1)m\) generators of \(J\).

\begin{proposition}[Normalized colength of the finite local condition]\label{prop:local-row-colength}
Let \(C_F\) be the cokernel on integral finite cohomology
of the map of Definition~\ref{hyp:first-level-row}. With
\(J\) the span of the exceptional classes constructed from
Proposition~\ref{prop:row-exceptional-coordinates}, the
normalized length satisfies
\[
 \frac{\operatorname{length}_O C_F}{\operatorname{length}_O(O/pO)}
       \le 2(D+1)rY+D(e+1)m\,18h .
\]
For \(p\ge3\), \(h\ge1\) and \(e\in\{1,2\}\), the common exponent
\(Y\) below is at most \(80(eh/p+eS+e\log_p h+1+\eta)\).
The additional \(18h\) allowance occurs on at most \(D(e+1)m\)
exceptional directions.
\end{proposition}

\begin{proof}
\proofstep{Lifting modulo the fixed exceptional classes.}
Let \((w,z)\) be a target Herr pair with \(w\in E^{-H}P_c\)
invariant to order \(M\) and \(z\in X^MN\). In the model
\(\beta(w)\), apply \(P(\gamma)\) and use the splitting into
main and residual coordinates. Expand the residual modulo
\(XN^{\rm mod}\) in the chosen basis of \(I_B\).
Return from the model to the Wach lattice using \(\alpha\). The integral multiplier and the cost of the two comparison maps yield
\[
 \frac{p^\kappa}{P(1)}\,P(1)p^{2\sigma}w=p^{2\sigma+\kappa}w.
\]

At exponent \(y_1\) every main tensor on its right lifts from
the source modulo \(XN_c\), by Section~\ref{par:row-mixed}.
The residual terms are integral combinations of the fixed
representatives. Multiply by \(p^\sigma\), raise all these
invariants to precision \(M\) and complete the source pair.
The difference of the first coordinates is invariant to order \(M\)
and lies in \(XN_c\). It lies in \(X^MN_c\) by uniqueness of
invariant lifting. Naturality and the deep coboundary criterion now give
\[
              p^{\,3\sigma+\kappa+y_1}[(w,z)]\in J
                       \quad\hbox{in } C_F .
\]
In particular the precision raising costs one \(\sigma\) for the
entire relation, irrespective of its number of terms.

\proofstep{Reduction of a finite class to the tested form.}
It remains to put a general finite target class into this form.
Its extension of the unit by \(T_2\) is crystalline over the fraction
field. Lift \(p^\rho\) times the unit generator in each component
to \(x_c\) in the extension Wach lattice. Its cyclotomic difference
lies in \(XN_c(T_2)\) since the quotient generator is fixed.
Here one needs the kernel intersection: the integral étale lattice
sequence is exact and the subrepresentation has the induced Wach
lattice. The same intersection holds in the Frobenius span, as the
coefficient ring is free over its Frobenius image. In the bounded
pole module it follows after multiplying by \(E^H\).
The graded correction produces \(x'_c\) with
\[
 x'_c-p^\sigma x_c\in XN_c(T_2),\qquad (\gamma-1)x'_c\in X^M N_c(T_2).
\]

The correction uses the graded scalar calculation on that kernel.
Its projection is still \(p^{\rho+\sigma}\) times the unit.
Consequently
\[
                        w_c=\varphi x'_{c-1}-x'_c
\]
belongs to the kernel in \(E^{-H}P_c\) and is invariant to order
\(M\), since \(\gamma\) commutes with \(\varphi\) and \(\varphi\)
preserves \(X^MN\). Choose the unique \(z_c\in X^MN\) with
\((\varphi-1)z_c=(\gamma-1)w_c\). On summing in \(c\), uniqueness
gives \(\sum z_c=(\gamma-1)\sum x'_c\), since both sides are deep
and have the same Frobenius difference. Thus
\(\sum(w_c,z_c)\) represents \(p^{\rho+\sigma}\) times the original
class by the connecting map, allowing its fixed unit. We have proved
\[
       p^Y C_F\subset J,\qquad
       Y=2\rho+4\sigma+\kappa+\max(2\sigma,y_c).
       \tag{R0}\label{eq:tag-r0}
\]

\proofstep{Bounding the exceptional submodule.}
A uniform exponent on \(J\) follows from the coarser full-level lift.
Over \(K_0(\mu_p)\) the two height bounds are \(h,h\), and the jumps
of sources and kernels remain in \([-h,h]\). Restriction and
corestriction cost the degree \(s\), a unit in \(O\).
In the full-level computation of Section~\ref{subsec:ramified-lift}
the two Wach losses are at most \(4h+v_p((4h)!)\) each.

The
kernel correction has loss at most
\[
 h+v_p(h!)+\sum_{j\in\mathcal J}\bigl(1+v_p(j)\bigr),
 \qquad
 \mathcal J\subset[-h,2h]\cap\mathbb Z,\quad 0\notin\mathcal J .
\]
Indeed the shift of the jump set uses degrees \(0,\ldots,h\) and
the positive degrees for the trivial component are at most \(h\);
zero eigenvalue differences are discarded in the rational coboundary
test. Even using the full set of nonzero signed integers in this
interval gives valuation sum at most \(v_p(h!)+v_p((2h)!)\).
Since \(p\ge3\), \(v_p(j!)\le j/2\). The total exponent is at most
\(18h\). Thus \(p^{18h}\) kills \(J\).

For clarity about the other rank factor, the rational Bloch--Kato
dimension formula bounds \(\dim H^1_f(K,V_2)\) by
\[
 \dim H^0(K,V_2)+\dim(D_{\rm dR}(V_2)/{\rm Fil}^0) \le (D+1)r.
\]
For a uniformizer \(\varpi\) of \(O\), the
cohomology sequence of multiplication by \(\varpi\) exhibits
\(H^1(K,T_2)[\varpi]\) as a quotient of \(H^0(K,T_2/\varpi T_2)\).
Its residue dimension is at most \(r\). Local integral cohomology
is finite type by local Tate finiteness. Its torsion and the
torsion of its finite-condition submodule therefore need at most
\(r\) generators by the structure theorem over a DVR.
Hence \(2(D+1)r\) generators suffice for the finite-condition
module and its quotients. Taking lengths first in \(C_F/J\) and
then in \(J\) gives
\[
 \frac{\operatorname{length}_O C_F}{\operatorname{length}_O(O/pO)}
     \le 2(D+1)rY+D(e+1)m\,18h .
\]
In particular the loss proportional to \(h\) per generator only
occurs on the \(D(e+1)m\) exceptional directions.

\proofstep{The common scalar exponent.}
For the numerical comparison put \(v=H/s\). Then
\(v\le 3eh/(2p)+eS/2+1\), as \(p\ge3\) and \(e\le p-1\).
The factorial bound gives \(\rho\le6v+3/2\) and \(\sigma\le6v\).
Therefore
\[
 Y\le48v+3+e(3S/2+\log_p H)+\eta+3/2+\log_p H .
\]

For \(h\ge1\), \(h+S+1\le h p^{S+1}\), so this is at most
\(80(eh/p+eS+e\log_p h+1+\eta)\) for \(e=1,2\).
The host always has \(h\ge1\). Substitution gives the ramified
estimate above and the same formula with \(e=1\) on an unramified
one-sign component. All choices precede the finite target class
apart from its own lift and expansion. The constants are independent
of the number of coefficient places, their ramification, the number
of Frobenius components, the source rank, and the outer jumps.
\label{cor:row-colength-1}
\end{proof}


\begin{proposition}[Weighted cost at the two ramified auxiliary primes]\label{prop:auxiliary-prime-row-length}
At \(\ell=q\) or \(d\), use the local field
\(E_w=F_w(\sqrt{-\ell})\) of \eqref{eq:host-auxiliary-completion}
for each \(w\mid\ell\), and
use the signed jump windows at each individual embedding
from Section~\ref{subsec:ramified-auxiliary-row}.
The local bound \eqref{eq:tag-r4}, after allowing the separate scalar
cohomology and class number budgets, is \(O(fmb)\).
Writing \(c_\ell\) for this local cost, the combined bound on the selected interval is
\[
 q,d\le\ell_*\quad\text{eventually},\qquad \sum_{\ell\in\{q,d\}}(\log\ell)c_\ell\ll fmb\log b.
\]
\end{proposition}

\begin{proof}
In this comparison put \(D_w=[E_w:\mathbb Q_\ell]\),
\[
 X_\ell=4b/\ell+2S_c+2\log_\ell(2b)+1+\eta_c+t,\qquad
 Y_b=2b+1.
\]
On the original row lattices the integral denominator
bound for a crystalline source and quotient with
the preceding jump conditions gives
\[
 L_w\le 160\bigl((D_w+1)r X_\ell+3D_wmY_b\bigr)+1 .
 \tag{R1}\label{eq:tag-r1}
\]
Here the local field is admitted at full level one,
with its coefficient field containing the embeddings
of \(F_w(\mu_\ell)\), and \(e(E_w/F_w)=2\).
The raw bound uses \(h_0=\lfloor2b\rfloor\),
\(S=\lfloor S_c\rfloor\) and
\(\lfloor\eta_c\rfloor+t\) in place of \(2b,S_c,\eta_c+t\).
Indeed the first two terms in \(X_\ell\) bound
\(2h_0/\ell+2S\), and its logarithmic term bounds
\(2\log_\ell h_0\). Rounding the resulting row bound up
adds at most one.

In the character rows again put \(h_0=\lfloor2b\rfloor\)
and write
\[
 \lambda=\lfloor\log_\ell(2h_0)\rfloor.
\]

The distinct jump set \(U\) of a source for a character
row lies in the union of the labeled jump sets. It
has at most \(D_wR\le2fR\) elements and lies in
\([-h_0,h_0-1]\). In the character lifting argument
at depth \(M=3+n_0\), take \(n_0\) least with
\[
 2h_0\le(\ell-1)\ell^{n_0},
\]
so \(M\le4+\lambda\).
The base change to \(F_w(\mu_\ell)\) costs no
\(\ell\)-power, as its relative degree is
\((\ell-1)/2\). The relative degree from
the first full level to depth \(M\) has valuation
\(M-1\).

For a twist exponent \(0\le e<t\), each nonzero
difference \(a-i\), with \(a\in U\) and
\(0\le i\le e\), has \(v_\ell(a-i)\le\lambda\).
The distinct-jump sum in each degree is bounded by
\(2fR(4+2\lambda)\). The exponent of the
truncation, scalar Bockstein and integral dual
exponential of Section~\ref{subsec:ramified-character}
is consequently at most
\[
 c_B=6+\lambda+7t+2tfR(4+2\lambda).
\]
Here \(6+\lambda\) allows the depth term and the
constant scalar terms, and \(7t\) allows the
per-degree and scalar terms linear in \(t\).
Since \(\lambda\le4\log b\) for \(\ell\ge3\), this gives
\[
                  c_B\le60tfR\log b .
 \tag{R2}\label{eq:tag-r2}
\]
The application is to the cocycle of the chosen
integral lift of the trivial quotient of a crystalline
extension. The character comparison moves its image
into finite cohomology. The rank-one local length
estimate then costs \((D_w+2)c_B\).
\label{cor:special-prime-lengths-9}

With no columns the global determinant bound of the
previous section has no column terms, and its row
length is exactly \(L_w\). The character row cost is
at most \(t\sum_{w\mid\ell}(D_w+2)c_B\).
The cyclotomic mark has depth one and \(\ell>t\).
Thus its normalized valuation is \(t-1\), giving a
determinant allowance at most \(4Rt\). The tame
character allowance is at most \(\ell^{-2}\). Together
with the signature length and the strict and unramified
global kernels this gives
\[
 v_{cc}(i,\pi)\le
 \sum_{w\mid\ell} L_w+t\sum_{w\mid\ell}(D_w+2)c_B
   +4Rt+\ell^{-2}+s_\pi+V_\ell+v_\ell(h_E).
 \tag{R3}\label{eq:tag-r3}
\]
Every \(D_w\ge2\) and \(\sum_wD_w=2f\).

In particular the two row sums are bounded by
\[
 160(3fr X_\ell+6fm Y_b)+f,\qquad
 240 f^2t^2R\log b
\]
respectively. The additional \(4Rt\) and \(f\)
are bounded by \(20 f^2t^2R\log b\).
For \(\ell\ge3\) and \(b\ge6\), each coefficient
in \(3X_\ell\) is at most \(13\) times its
corresponding coefficient in
\(b/\ell+S_c+\eta_c+t+\log b\), absorbing the constant
into the logarithm. Also \(6Y_b\le13b\) once \(b\ge6\).
Hence at this
threshold \eqref{eq:tag-r3} yields
\[
 \begin{split}
 v_{cc}(i,\pi)\le{}&
 2080 f\bigl(r(b/\ell+S_c+\eta_c+t+\log b)+mb\bigr)
     +260 f^2t^2R\log b\\
 &+2\ell^{-2}+V_\ell+v_\ell(h_E).
 \end{split}
 \tag{R4}\label{eq:tag-r4}
\]
For the second \(\ell^{-2}\) we used the signature
bound. This supplies the entry at the ramified quadratic prime on the same point
as its \(v_{cc}\) in the cut lattice.

\proofstep{Use the prime range to absorb the non-leading two terms in \eqref{eq:tag-r4}.}
For the selected auxiliary prime pair,
\[
 \log q,\log d=o(\log b)
\]
and both primes exceed
every fixed power of \(x\), uniformly in the scale
interval of that pair. In particular the condition
\(\ell>x^{2C_t+5}\) at the pointwise ramified caller
holds for the fixed rank exponent \(C_t\).
The ranks grow polynomially,
and \(\eta_c,S_c,t\) are at most
\(x^{O(1)}\exp(x^{\epsilon/4})\).
Thus the two terms of \eqref{eq:tag-r4} other than \(2080fmb\)
and the common budgets are \(o(fmb)\), uniformly
over both primes and all coefficient localizations.

In particular the table cost is at most \(2Cfmb\)
at a fixed large-scale threshold.
Also
\[
 q,d\le\ell_*,
\]
since the denominator of
\(\ell_*\) is \(x^{O(1)}\exp(O(x^{\epsilon/4}))\)
and its numerator is \(mb\). Summing with
\(\log q,\log d\le\log b\) gives \(O(fmb\log b)\),
which fits the stated \(O(fmb\log^2 b)\) budget.
\label{cor:special-prime-lengths-10}
\end{proof}

